% ============================================================================
% DRAFT SECTIONS FOR NEW ANALYSES -- SPLICED 2026-08-03. ARCHIVE COPY ONLY.
% Edit paper_draft_working.tex, not this file. Kept for provenance.
%
% Splice record (2026-08-03):
%   ORDER  -> methods half is Sec 3.7 \subsection{Alternative ordering}
%             (sec:alternative-order); results half is Sec 4.12
%             \subsection{Alternative Ordering} (sec:res-order). Numbers were
%             refreshed to the 5-run result before splicing; see the block
%             comment below.
%   B4/B5 realised-weight -> NOT spliced. The main draft already carries
%             tab:swing_weights (Sec 2.2) for realised swing and the newly
%             added tab:zero_variance (app:weights) for the discrimination
%             rates, so tab:realised_weight would have duplicated both.
%   B5 sensitivity table -> spliced as tab:sensitivity_by_model, but rebuilt
%             PER MODEL rather than as a four-model mean, per user instruction
%             not to pool across models. Top-1 for the same cells went to
%             supplementary tab:sensitivity_top1. The +/-0.05 rows are deleted
%             from the main paper and the supplementary.
%   B6     -> Discussion, Interpretation of Results.
%   B7     -> app:weights prose, with two corrections to the block as drafted:
%             (a) the "three of 195 top-1 rankings" figure was unverifiable
%                 from the repo and was dropped; the 285-scope justification
%                 now rests on calibration consistency, which was verified
%                 (HVAC active-cost p5/p95 over all 285 = 0.3820/3.2920,
%                 matching the calculator's $0.38--$3.29).
%             (b) "implied weights collapse to corner solutions, placing the
%                 entire weight on Comfort or Practicality" is false for
%                 Shower, which returns (0.211, 0.789, 0, 0). Corrected.
%
% HISTORY: earlier drafts of this file carried a position-bias Results and
% Discussion block built on an A_E arm in which retrieval was silently dead
% (chroma_collection was None; "Retrieved 0 similar scenarios" for all 585 calls
% per model). That arm measured exemplar removal, not alternative order, and
% those blocks were removed. The A_D arm from the same runs is valid, because
% A_D never retrieves, and it is the empirical support cited in the ORDER block.
%
% CITATIONS: wang2023positionbias and zheng2023judging were added to
% paper/cas-refs.bib on 2026-08-03. No longer blocking.
% ============================================================================


% ============================================================================
% B1-B3 REPLACEMENT -- ALTERNATIVE-ORDER JUSTIFICATION
% Destination: Methods, immediately after the prompt design description.
% Supporting data: Analysis/PositionBias/position_bias_summary_all_models.xlsx
%                  (A_D arm, 195 scenarios x 4 models, valid)
% ============================================================================

\subsection{Alternative ordering}
\label{sec:alternative-order}

Language models weight response options by position as well as by content
\cite{wang2023positionbias,zheng2023judging}, so a fixed alternative order is a
candidate confound. The exposure depends on how each architecture elicits
scores, and it differs sharply across the three.

$\mathcal{A}_{\text{D}}$ and $\mathcal{A}_{\text{E}}$ score one alternative per
API call. Each call names the alternative under evaluation, states the scenario
context, and lists the remaining two alternatives on a single line for context.
The calls are independent and carry no conversation state, so the sequence in
which they are issued cannot affect any individual response. Reversing the
alternative order therefore changes exactly one thing the model can observe:
the order of a two-element list mentioned in passing. Scoring a setpoint of
76\,\textdegree F, for example, the prompt moves from ``Other alternatives
available for this decision: [72, 80]'' to ``[80, 72]'' while the alternative
being scored, the scenario context, and the scoring instructions are
byte-identical. Neither architecture ever presents the three alternatives as a
ranked list to choose among, which is the presentation format the position-bias
results above concern. The mechanism those results describe has no route into
this design.

We verified this rather than assuming it. Rescoring all 195 test scenarios with
the alternative order reversed, across all four models, left
$\mathcal{A}_{\text{D}}$ unchanged: the reversed arm was less accurate in one of
four models, and no model separated from its own five-run variation (exact
McNemar on top-1 correctness, $p \geq 0.807$ against every shipped run). Because
the ranking tie-break resolves exact ties by input order, reversal can flip a
fully tied choice set on its own; that artifact biases the comparison toward
detecting an effect, so the observed null is conservative.

$\mathcal{A}_{\text{H}}$ is not covered by that argument. Its extraction call
passes the whole scenario record, including all three alternatives as a numbered
sequence, so it is the one architecture whose prompt matches the format the
position-bias literature describes. Any effect would be indirect, acting through
the estimated physical parameters rather than through a score, because the
calculator and not the model produces the ranking. We therefore tested it
directly.

We re-ran $\mathcal{A}_{\text{H}}$'s extraction over all 195 test scenarios with
the alternative values reversed in the prompt, five times per model, and paired
it with three same-session runs in the shipped order. The control arm is
necessary because the benchmark runs were collected days earlier: a difference
between a reversed arm and an older shipped arm is equally consistent with an
ordering effect and with provider-side drift. Extracted parameters were handed
to the calculator with the alternatives in canonical order in every arm, so the
ranking tie-break cannot manufacture a difference between them. We test whether
reversed runs are distinguishable from shipped-order runs by relabelling which
runs carry the reversed label across all $\binom{13}{5}$ assignments, an exact
permutation test that uses every scenario rather than only the few on which two
runs disagree.

Ordering changed what the models extracted, and did not change what they
decided. Two of four models separated at the parameter level (Qwen and Gemini,
separation $0.054$ and $0.047$, $p = 0.0008$, Holm-adjusted $p = 0.006$ across
the eight tests), while no model separated on the top-1 alternative
($p \geq 0.24$). Both significant $p$-values sit at the exact permutation floor
of $1/1287$: the observed labelling was the most extreme of all assignments, so
the true $p$ is bounded above by that figure rather than estimated at it. The
separation is a shift in which values the model returns, not a loss of accuracy:
against the reference ranking, the reversed and control arms differ by less than
$0.004$ Kendall's $\tau$, and the sign of that difference is not consistent
across models. Listing the alternatives in a
different order moves Gemini's estimate of a dwelling's insulation and equipment
age, quantities that have nothing to do with the order the setpoints are printed
in. The deterministic calculator absorbs that perturbation before it reaches the
ranking, which is the behaviour the architecture is designed to produce and
which the two model-side null results alone would not have demonstrated.

One incidental observation bears on any benchmark that compares arms collected
at different times. For DeepSeek, run-to-run agreement on extracted parameters
was $0.42$ among the shipped runs and $0.64$ and $0.67$ among the two arms
collected in the later session, with no manipulation distinguishing the latter
two. Had we compared a
reversed arm against the shipped runs without a contemporaneous control, that
gap would have been readable as an effect of the manipulation.

% NUMBERS CURRENT as of 2026-08-03, 5 reversed + 5 shipped + 3 control runs per
% model (Analysis/Hybrid_Ablation/hybrid_order_reversal.xlsx, sheet "summary",
% pooled_* columns; C(13,5) = 1287 relabelings):
%   pooled_param_sep / pooled_param_p:
%     qwen  +0.0543 / 0.000777   gemini   +0.0471 / 0.000777
%     gptoss +0.0041 / 0.302     deepseek -0.0158 / 0.450
%   pooled_choice_p: gemini 0.235 (min), qwen 0.679, gptoss 0.838, deepseek 0.953
%   Holm over 8 tests: qwen param 0.0062, gemini param 0.0062, all others 1.000
%   tau order_control vs order_reversed:
%     deepseek 0.8974/0.9009  gemini 0.9293/0.9262
%     gptoss   0.8988/0.8970  qwen   0.8895/0.8913   (max |gap| 0.0035)
%   deepseek param identity: within_shipped 0.4230, within_control 0.6393,
%     within_reversed 0.6703


% ============================================================================
% B4  METHODS -- WEIGHT REALISATION DIAGNOSTICS
% Script: Miscellaneous Scripts/WeightDiagnostics.py
% Output: Scoring Logic and Documentation/method/weight_discrimination.xlsx
%         Scoring Logic and Documentation/method/cost_env_collinearity.xlsx
% ============================================================================

\subsubsection{Realised weight}
\label{sec:realised-weight}

A criterion weight only affects a ranking to the extent that the criterion
varies across the alternatives being ranked. We therefore report, alongside the
nominal weights, the mean within-scenario range of each criterion and the
product of the two, which we call realised swing. Normalising realised swing to
sum to one puts it on the same scale as the nominal vector, so the two are
directly comparable.


% ============================================================================
% B5  RESULTS -- WEIGHT REALISATION AND SENSITIVITY
% Regenerate: python "Miscellaneous Scripts/WeightDiagnostics.py"
%             python "Miscellaneous Scripts/SensitivityAnalysis.py" --all-models
% ============================================================================

\FloatBarrier

\begin{table}[htbp]
\centering
\begin{threeparttable}
\footnotesize
\caption{Nominal weight against realised swing share, by decision type.
Realised swing is the nominal weight times the criterion's mean within-scenario
range, normalised to sum to one. ``No discrim.'' is the percentage of scenarios
in which the criterion takes the same value for all three alternatives.}
\label{tab:realised_weight}
\begin{tabular}{@{}llrrr@{}}
\toprule
Type & Criterion & Nominal & Realised & No discrim. (\%) \\
\midrule
HVAC & Energy cost    & 0.300 & 0.172 & 10.5 \\
     & Environmental  & 0.350 & 0.198 & 10.5 \\
     & Comfort        & 0.200 & 0.538 & 0.0  \\
     & Practicality   & 0.150 & 0.092 & 44.8 \\
\midrule
Appliance & Energy cost   & 0.300 & 0.247 & 28.0 \\
          & Environmental & 0.350 & 0.029 & 42.0 \\
          & Comfort       & 0.200 & 0.368 & 0.0  \\
          & Practicality  & 0.150 & 0.356 & 0.0  \\
\midrule
Shower & Energy cost   & 0.300 & 0.301 & 1.2 \\
       & Environmental & 0.350 & 0.357 & 1.2 \\
       & Comfort       & 0.200 & 0.192 & 0.0 \\
       & Practicality  & 0.150 & 0.149 & 0.0 \\
\bottomrule
\end{tabular}
\end{threeparttable}
\end{table}

Table~\ref{tab:realised_weight} shows the design weights are realised as
specified in one decision type and not in the other two. For Shower the
realised shares reproduce the nominal vector to within 0.008 across all four
criteria. For HVAC, Comfort realises 0.538 against a nominal 0.200 while
Practicality contributes almost nothing and fails to separate the alternatives
in 44.8\% of scenarios. For Appliance, Environmental realises 0.029 against the
largest nominal weight in the design, 0.350, because 42\% of Appliance
scenarios contain no environmental difference between the alternatives at all.

This is the same structure the objective weight methods report from a different
direction. MEREC's zero-variance counts match these discrimination rates
exactly, and MEREC's per-type vectors load onto Comfort for HVAC (0.663) and
away from Environmental for Appliance (0.044), which is what a method that
weights by within-scenario discrimination must do given
Table~\ref{tab:realised_weight}.

% --- Replacement sensitivity table. The +/-0.05 rows are removed: they span
% --- 0.179 to 0.225 on the between-architecture gap while the data-derived
% --- arms reach 0.090, so they cannot establish weight robustness on their own.

\begin{table}[htbp]
\centering
\begin{threeparttable}
\footnotesize
\caption{Mean Kendall's $\tau$ across the four evaluated models under
data-derived weight vectors, with the between-architecture gap.}
\label{tab:sensitivity_pooled}
\begin{tabular}{@{}lrrrrr@{}}
\toprule
Weight vector & $\mathcal{A}_{\text{D}}$ & $\mathcal{A}_{\text{E}}$ & $\mathcal{A}_{\text{H}}$ & $\mathcal{A}_{\text{E}}-\mathcal{A}_{\text{D}}$ & $\mathcal{A}_{\text{H}}-\mathcal{A}_{\text{E}}$ \\
\midrule
Design (baseline)   & 0.095 & 0.315 & 0.902 & 0.220 & 0.587 \\
Equal               & 0.139 & 0.345 & 0.913 & 0.205 & 0.568 \\
\midrule
Entropy pooled      & 0.109 & 0.305 & 0.874 & 0.197 & 0.569 \\
Entropy per-type    & 0.074 & 0.323 & 0.902 & 0.250 & 0.579 \\
\midrule
MEREC pooled        & 0.327 & 0.455 & 0.942 & 0.127 & 0.487 \\
MEREC per-type      & 0.335 & 0.463 & 0.936 & 0.128 & 0.473 \\
MEREC HVAC vector   & 0.403 & 0.493 & 0.968 & 0.090 & 0.475 \\
\bottomrule
\end{tabular}
\begin{tablenotes}
\footnotesize
\item Each cell is the arithmetic mean of the four per-model values. Both the
reference ranking and the architecture ranking are recomputed with each vector,
because a MAVT reference is defined by its weights.
\end{tablenotes}
\end{threeparttable}
\end{table}

The architecture ordering
$\mathcal{A}_{\text{H}} > \mathcal{A}_{\text{E}} > \mathcal{A}_{\text{D}}$ holds
under every weight vector tested, including MEREC's HVAC vector, which assigns
Comfort 0.663 against the design's 0.200. The magnitude of the
$\mathcal{A}_{\text{E}}$ advantage, however, is weight-dependent in a way the
design vector conceals. Under MEREC weights $\mathcal{A}_{\text{D}}$'s $\tau$
rises from 0.095 to 0.403 while $\mathcal{A}_{\text{E}}$'s rises only from
0.315 to 0.493, so the gap between them falls from 0.220 to 0.090.


% ============================================================================
% B6  DISCUSSION -- WEIGHTS
% ============================================================================

\FloatBarrier

Weighting toward the criterion that separates alternatives closes most of the
distance between direct prompting and example-guided scoring. MEREC vectors
load onto Comfort; entropy vectors load onto Energy Cost and Environmental
Impact, as the design does, and produce results indistinguishable from the
design vector. The architectures are therefore not uniformly better or worse at
multi-criteria judgement. They differ most where the criteria are hardest for a
language model to estimate, which in this corpus is cost and emissions, and
least on Comfort.

That reading has a direct consequence for what
$\mathcal{A}_{\text{E}}$ contributes. Its advantage over
$\mathcal{A}_{\text{D}}$ is concentrated in estimating quantities that follow
from tariff structure and grid emissions factors rather than in ranking
alternatives per se. The design weights place 0.65 on those two criteria, which
makes the reported benchmark the least favourable weighting for
$\mathcal{A}_{\text{D}}$ among those tested, and the architecture ordering
survives it.

The Appliance case also shows the cost of the criterion definitions chosen to
avoid collinearity. Holding emissions to a marginal-factor definition on a
7am--11pm window while cost follows utility-specific peak windows leaves 42\%
of Appliance scenarios with no environmental difference between alternatives,
so the criterion carrying the largest nominal weight contributes least to the
ranking there.


% ============================================================================
% B7  APPENDIX -- CORRECTED WEIGHT-ROBUSTNESS PROSE
% Replaces the paragraph currently ending "Entropy and MEREC therefore provide
% the informative objective comparisons here". Corrects three errors:
%   (1) MEREC does NOT assign above-average weight to Environmental (0.131).
%   (2) ImpliedWeights.py uses constrained least squares, not multinomial
%       logistic regression; there is no likelihood.
%   (3) The corpus scope was never stated.
% ============================================================================

Entropy and MEREC weights are computed over the full 285-scenario corpus, the
same corpus the value functions are calibrated against; restricting the
percentile bounds to the 195 evaluated scenarios changes three of 195 reference
top-1 rankings.

Implied weights are derived by constrained least squares on pairwise
preferences between alternatives, minimising squared error subject to
non-negativity and a unit sum. When a criterion takes the same value for all
alternatives in a scenario it contributes no discriminatory information, and
the fit concentrates on whichever criterion does vary.

The three objective methods disagree with the design in different directions,
and the disagreement is informative rather than corrective. Entropy weights
support the Environmental priority, assigning it 0.371 against the design's
0.350. MEREC does not: it assigns Environmental 0.131, the lowest of its four
criteria, and Comfort 0.404 against the design's 0.200. Implied weights
collapse to corner solutions, placing the entire weight on Comfort or
Practicality, and the Appliance fit reports a negative pairwise $R^2$
($-0.211$), so they are reported here as a discrimination diagnostic rather
than as a candidate weighting.

The two disagreements share one cause. Entropy measures spread across the whole
corpus, where cost and emissions vary widely between scenarios. MEREC and the
implied fit measure spread within a choice set, where they often do not. The
design weights follow the entropy reading, and the sensitivity analysis
(Table~\ref{tab:sensitivity_pooled}) reports what happens under the MEREC
reading instead.
